% Part: computability
% Chapter: computability-theory
% Section: ce-closed-cup-cap

\documentclass[../../../include/open-logic-section]{subfiles}

\begin{document}

\olfileid{cmp}{thy}{clo}

\olsection[اتحاد المجموعات القابلة للتعداد حسابيًا وتقاطعها]{المجموعات
  القابلة للتعداد حسابيًا منغلقة تحت الاتحاد والتقاطع}

للمجموعات القابلة للتعداد حسابيًا خواص انغلاق،
منها ما تقرره المبرهنة الآتية.

\begin{thm}
متى كانت $\OLId{latin-looped}{A}$ و$\OLId{latin-looped}{B}$ قابلتين للتعداد حسابيًا، كان كل من
$\OLId{latin-looped}{A}\cap \OLId{latin-looped}{B}$ و$\OLId{latin-looped}{A}\cup \OLId{latin-looped}{B}$ كذلك.
\end{thm}

\begin{proof}
نستعمل \olref[eqc]{thm:ce-equiv} لننتقل بين توصيفات المجموعات
القابلة للتعداد حسابيًا، ونورد لذلك براهين مختلفة. وإذا خلت إحدى المجموعتين فالنتيجة ظاهرة؛ فلنفرضهما غير خاليتين في البرهانين الأولين.

أما البرهان الأول، فلنأخذ دالة قابلة للحساب~$\OLId{latin-ordinary}{f}$ تعدد $\OLId{latin-looped}{A}$، ودالة
قابلة للحساب~$\OLId{latin-ordinary}{g}$ تعدد $\OLId{latin-looped}{B}$، ونضع
\begin{align*}
\OLId{latin-ordinary}{h}(\OLId{latin-ordinary}{x}) & = \umin{\OLId{latin-ordinary}{y}}{(\OLId{latin-ordinary}{f}(\OLId{latin-ordinary}{y}) = \OLId{latin-ordinary}{x} \lor \OLId{latin-ordinary}{g}(\OLId{latin-ordinary}{y}) = \OLId{latin-ordinary}{x})} \text{، و}\\
\OLId{latin-ordinary}{j}(\OLId{latin-ordinary}{x}) & = \umin{\OLId{latin-ordinary}{y}}{(\OLId{latin-ordinary}{f}((\OLId{latin-ordinary}{y})_\OLMathNumeral{0}) = \OLId{latin-ordinary}{x} \land \OLId{latin-ordinary}{g}((\OLId{latin-ordinary}{y})_\OLMathNumeral{1}) = \OLId{latin-ordinary}{x})}.
\end{align*}
عندئذ يكون $\OLId{latin-looped}{A}\cup \OLId{latin-looped}{B}$ مجال $\OLId{latin-ordinary}{h}$، ويكون $\OLId{latin-looped}{A}\cap \OLId{latin-looped}{B}$ مجال~$\OLId{latin-ordinary}{j}$.

\begin{explain}
وتفسير ذلك بالحساب: إذا كان لنا إجراءان لتعداد $\OLId{latin-looped}{A}$ و
$\OLId{latin-looped}{B}$، شبه قررنا انتماء $\OLId{latin-ordinary}{x}$ إلى $\OLId{latin-looped}{A}\cup \OLId{latin-looped}{B}$ بطلب
$\OLId{latin-ordinary}{x}$ في أي من التعدادين؛ وشبه قررنا انتماء $\OLId{latin-ordinary}{x}$ إلى
$\OLId{latin-looped}{A}\cap \OLId{latin-looped}{B}$ بطلب $\OLId{latin-ordinary}{x}$ في كليهما معًا.
\end{explain}

وأما البرهان الثاني، فنبقي~$\OLId{latin-ordinary}{f}$ معددة لـ$\OLId{latin-looped}{A}$ و~$\OLId{latin-ordinary}{g}$ معددة لـ$\OLId{latin-looped}{B}$،
ونضع
\[
\OLId{latin-ordinary}{k}(\OLId{latin-ordinary}{x}) = \begin{cases}
\OLId{latin-ordinary}{f}(\OLId{latin-ordinary}{x}/\OLMathNumeral{2}) & \text{إذا كان $x$ زوجيًا} \\
\OLId{latin-ordinary}{g}((\OLId{latin-ordinary}{x}-\OLMathNumeral{1})/\OLMathNumeral{2}) & \text{إذا كان $x$ فرديًا.}
\end{cases}
\]
فتعدد $\OLId{latin-ordinary}{k}$ المجموعة $\OLId{latin-looped}{A}\cup \OLId{latin-looped}{B}$، لأن $\OLId{latin-ordinary}{k}$ تتعاقب فيها قيمتا
التعدادين $\OLId{latin-ordinary}{f}$ و~$\OLId{latin-ordinary}{g}$. وأما $\OLId{latin-looped}{A}\cap \OLId{latin-looped}{B}$ فتحتاج إلى مزيد نظر: إن
كانت $\OLId{latin-looped}{A}\cap \OLId{latin-looped}{B}$ خالية، كانت قابلة للتعداد حسابيًا على نحو بديهي. وإلا
فليكن $\OLId{latin-ordinary}{c}$ أي عنصر من $\OLId{latin-looped}{A}\cap \OLId{latin-looped}{B}$، وعرف $\OLId{latin-ordinary}{l}$ بـ
\[
\OLId{latin-ordinary}{l}(\OLId{latin-ordinary}{x}) = \begin{cases}
\OLId{latin-ordinary}{f}((\OLId{latin-ordinary}{x})_\OLMathNumeral{0}) & \text{إذا كان $f((x)_0) = g((x)_1)$} \\
\OLId{latin-ordinary}{c} & \text{في غير ذلك.}
\end{cases}
\]
فالدالة $\OLId{latin-ordinary}{l}$ تستعرض أزواج عناصر تعدادي $\OLId{latin-ordinary}{f}$ و$\OLId{latin-ordinary}{g}$، وتخرج
القيمة المشتركة كلما اتفقتا؛ وإن لم تتفقا أخرجت $\OLId{latin-ordinary}{c}$ حتى لا تقف.

وأما البرهان الأخير، فلنأخذ $\OLId{latin-looped}{A}$ \emph{مجالًا} للدالة الجزئية القابلة للحساب $\OLId{latin-ordinary}{m}(\OLId{latin-ordinary}{x})$،
و$\OLId{latin-looped}{B}$ مجالًا للدالة الجزئية القابلة للحساب $\OLId{latin-ordinary}{n}(\OLId{latin-ordinary}{x})$. فتكون $\OLId{latin-looped}{A}\cap \OLId{latin-looped}{B}$ مجال الدالة
الجزئية $\OLId{latin-ordinary}{m}(\OLId{latin-ordinary}{x})+\OLId{latin-ordinary}{n}(\OLId{latin-ordinary}{x})$.

\begin{explain}
ووجهه أن $\OLId{latin-looped}{A}$ تجمع القيم التي تتوقف عندها $\OLId{latin-ordinary}{m}$،
و$\OLId{latin-looped}{B}$ تجمع القيم التي تتوقف عندها $\OLId{latin-ordinary}{n}$؛ فلا يكون $\OLId{latin-looped}{A}\cap \OLId{latin-looped}{B}$ إلا مجموعة القيم
التي يتوقف عليها الإجراءان جميعًا.
\end{explain}

وأما رد $\OLId{latin-looped}{A}\cup \OLId{latin-looped}{B}$ إلى مجموعة قيم توقف فيحتاج إلى
محاكاة $\OLId{latin-ordinary}{m}$ و$\OLId{latin-ordinary}{n}$ على التوازي. فنأخذ $\OLId{latin-ordinary}{d}$ فهرسًا لـ$\OLId{latin-ordinary}{m}$ و$\OLId{latin-ordinary}{e}$ فهرسًا
لـ$\OLId{latin-ordinary}{n}$، فيكون $\OLId{latin-ordinary}{m}=\cfind{\OLId{latin-ordinary}{d}}$ و$\OLId{latin-ordinary}{n}=\cfind{\OLId{latin-ordinary}{e}}$. عندئذ يكون $\OLId{latin-looped}{A}\cup \OLId{latin-looped}{B}$ مجال
الدالة
\[
\OLId{latin-ordinary}{p}(\OLId{latin-ordinary}{x}) = \umin{\OLId{latin-ordinary}{y}}{(\OLId{latin-looped}{T}(\OLId{latin-ordinary}{d},\OLId{latin-ordinary}{x},\OLId{latin-ordinary}{y}) \lor \OLId{latin-looped}{T}(\OLId{latin-ordinary}{e},\OLId{latin-ordinary}{x},\OLId{latin-ordinary}{y}))}.
\]
\begin{explain}
فالدالة $\OLId{latin-ordinary}{p}$ تطلب، عند الدخل $\OLId{latin-ordinary}{x}$، سجل حساب متوقف لـ$\OLId{latin-ordinary}{m}$ أو
لـ$\OLId{latin-ordinary}{n}$؛ ومتى عثرت على أحدهما توقفت.
\end{explain}
\end{proof}

\end{document}
